\documentclass{subfiles}
\begin{document}
    We begin by recalling the concepts of radical ideal. Let us also note that, 
        in the remaining part of these notes, we will be considering Noetherian rings 
        unless otherwise specified.
    
    \begin{definition*}
        Let \(R\) be a ring and let \(I \subseteq R\) be an ideal.
            The set 
            \[
                \sqrt{I} = \set*{x \in R}[x^{k} \in I, k \in \N]
            \]
    \end{definition*}

    \begin{proposition*}
        Given \(R\) a ring and \(I \subseteq R\) an ideal, 
            the set \(\sqrt{I}\) is an ideal and \(\sqrt{I} \supseteq I\).
    \end{proposition*}

    \begin{proof*}
        We have to prove that \(\sqrt{I}\) is closed under addition and multiplication.
        \begin{description}
            \item [Addition:] let \(x, y \in \sqrt{I}\). Then, \(\exists n, m \in \N\)
                such that \(x^{n}, y^{m} \in I\). Hence, we have that 
                \[
                    (x + y) \in \sqrt{I} \iff (x + y)^{m + n} \in I.
                \]
                Observe that
                \[
                    (x + y)^{m + n} = \sum_{k = 1}^{m + n} \binom{m + n}{k} x^{m + n - k}y^{k}.
                \]
                from which 
                \begin{enumerate}
                    \item if \(k \le m\), then \(x^{m + n - k}y^{k} = 
                        \underbrace{x^{n}}_{\in I} \underbrace{x^{m - k} y^{k}}_{\in R}\).
                    \item if \(m + 1 \le k \le m + n\), then \(x^{m + n - k}y^{k} = 
                        \underbrace{y^{m}}_{\in I} \underbrace{x^{m + n - k}y^{k - m}}_{\in R}\).
                \end{enumerate}
            Therefore, any term of the summation is in \(I\). Thus, \((x + y)^{m + n} \in I\).

            \item [Multiplication:] let \(x \in \sqrt{I}, a \in R\). We have that 
                \[
                    ax \in \sqrt{I} \iff (ax)^{n} \in I. 
                \]
                But 
                \[
                    (ax)^n = a^{n}x^{n} \in I,
                \]
                therefore \(ax \in \sqrt{I}\).\qed
        \end{description}
    \end{proof*}
    \clearpage

    \begin{proposition*}
        Given \(R\) a ring, and \(I \subseteq R\) an ideal, it holds that 
        \[
            I \text{ is radical} \iff \frac{R}{I} \text{ is reduced}.
        \]
    \end{proposition*}
    \begin{proof*}
        Trivial. Simply consider the epimorphism \(\pi : I \to \frac{R}{I}\).
            Note that \(\ker{\pi} = I\), from which the thesis follows easely.\qed 
    \end{proof*}

    \begin{definition*}[Affine algebraic set]
        \begin{enumerate}
            \item Let \(\F\) be a field and let \(\K = \F[x_{1}, x_{2}, \ldots, x_{n}]\).
                Let \(F(x_{1}, \ldots, x_{n}) \in \K\). 
                If \((a_{1}, \ldots, a_{n}) \in \F^{n}\) is such that \(F(a_{1}, \ldots, a_{n}) = 0\),
                we say that the point \(P = (a_{1}, \ldots, a_{n}) \in \mathbb{A}^{n}\) is a zero for \(F\).

            \item Let \(S \subseteq \K\) be a set of polynomials.
                The set 
                \[
                    V(S) = \set*{(a_{1}, \ldots, a_{n}) \in \F^{n}}[\forall F \in S, F(a_{1}, \ldots, a_{n}) = 0]
                \]
                is called affine algebraic set.
        \end{enumerate}
    \end{definition*}

    \begin{proposition*}
        Given \(S \subseteq \K = \F[x_{1}, x_{2}, \ldots, x_{n}]\) a set of polynomials,
            the sets \(V(S), V(I(S)), V(\set{G_{1}, \ldots, G_{n}})\) coincide.
    \end{proposition*}

    \begin{proof*}
        We prove that, if \(I(S)\) is generated by \(S\), then \(V(S) = V(I(S))\).

        Consider \(F \in I(S)\), we can write it as \(F = \sum_{i = 1}^{n} F_{i}a_{i}, F_{i} \in S\).
        \begin{description}
            \item [(\(\subseteq\))] If \(P \in V(S)\), then \(P\) is a zero for any polynomial
                in \(S\); that is, \(F(P) = 0\). Hence, \(P \in V(I(S))\).
            \item [(\(\subseteq\))] If \(P \in V(I(S))\), then \(F(P) = 0\) for any \(F \in I\).
                By hypothesis \(I\) is generated by \(S\), meaning \(S \subseteq I\).
                Hence, \(G(P) = 0, \forall G \in S\).\qed
        \end{description}
    \end{proof*}

    \begin{proposition*}
        Affine algebraic sets satisfy the properties of closed sets in a topology.
    \end{proposition*}

    \begin{proof*}
        What we have to prove is that, the union of finitely many affine sets,
            is an affine set; and likewise, that the intersection of finitely many 
            affine sets is a affine set.

        Let us consider the intersection first. Let \(V(S_{1}), \ldots, V(S_{k})\) be a finite 
            sequence of affine sets. Let 
            \[
                V(S) = \bigcap_{i = 1}^{k} V(S_{i})
            \]
            be the intersection of all \(V(S_{i})\). We prove that if \(P \in V(S)\), 
            then \(P \in V(S_{i}), \forall i\). 

        By the definition of affine set 
        \[
            V(S) = \set*{P \in \mathbb{A}^{n}}[f(P) = 0, \forall f \in S],
        \]
        meaning that 
        \[
            V(S) = \bigcap_{i = 1}^{k} V(S_{i}) = V \left\lparen \bigcup_{i = 1}^{k} S_{i} \right\rparen .
        \]

        For what regards the union, we can focus on the simple case of the union of two affine sets.
            Let \(V_{1} = V(F_{1}, \ldots, F_{k}), V_{2} = V(G_{1}, \ldots, G_{l})\) be affine sets. 
            We prove that \(V_{1} \cup V_{2} = V = V(F_{1}, \ldots, F_{k}, G_{1}, \ldots,G_{l})\).

        \begin{description}
            \item [\(\subseteq\):] Let \(P \in V_{1} \cup V_{2}\). 
                Then, \(P \in V_{1} \lor P \in V_{2}\). In any case \(F_{i}(P)G_{j}(P) = 0\) 
                for any \(i, j\). Meaning \(P \in V\).

            \item [\(\supseteq\):] Let \(P \in V\). Then \(F_{i}(P)G_{j}(P) = 0\) for any
                \(i, j\). Meaning that 
                \begin{itemize}
                    \item If \(F_{i}(P) = 0, P \in V_{1}\), or 
                    \item If \(F_{i}(0) \ne 0 \implies G_{j}(P) = 0, P \in V_{2}\).
                \end{itemize}
        \end{description}
        In any case \(P \in V_{1} \cup V_{2}\).\qed
    \end{proof*}

    \begin{definition*}[Zarinski topology]
        A topology of \(\mathbb{A}^{n}\) is called a Zarinski topology if its closed sets 
            are affine algebraic sets.

        If \(X \subseteq \mathbb{A}^{n}\) is a subset, its Zarinski topology is the one
            induced by the one of \(\mathbb{A}^{n}\). Its closed sets have the shape 
            \[
                X \cap S(\K), \K = \F[x_{1}, \ldots, x_{n}].
            \]
    \end{definition*}

    \begin{definition*}
        Let \(V \subseteq \mathbb{A}^{n}\). 
            The set 
            \[
                I(V) = \set*{f \in \K}[f(P) = 0, \forall P \in V]
            \]
            is the set of those polynomials that go to zero in all points of  \(V\).
    \end{definition*}

    \begin{remark*}
        It can be proved that \(I(V)\) is an ideal. Also, two maps can be defined.
        \subfile{../../extra/TikZ/Figure - Commutative diagram ideal-affine sets}
    \end{remark*}
    \clearpage

    \begin{proposition*}
        Let \(X \subseteq \mathbb{A}^{n}\), then 
        \begin{enumerate}
            \item \(X \subseteq V(I(X))\);
            \item \(X = V(I(X)) \iff X\) closed set;
            \item \(I(X)\) radical ideal;
            \item There exist \(\phi\) one-to-one:
                \subfile{../../extra/TikZ/Figure - Commutative diagram properties ideals}
        \end{enumerate}
    \end{proposition*}

    \begin{proof*}
        \begin{enumerate}
            \item Trivially true. In fact, if \(F(x_{1}, \ldots, x_{n}) \in I(X)\)
                this means that \(F(P) = 0, \forall P \in X\).
                Hence, a point \(P \in X\) is such that \(G(P) = 0, \forall G \in I(X)\).
                Thus, \(\forall P \in X, P \in V(I(X))\).

            \item Let us note that the left implication is true by definition.
                The other implication follows. Let \(X = V(S)\) closed affine set.
                By definition, \(S \subseteq I(X)\), which means \(X \subseteq V(I(X)) \subseteq V(S)\).
                But we've said that \(X = V(S)\); therefore, \(V(I(X)) \subseteq X \implies X = V(I(X))\).

            \item Let \(F^{n}(x_{1}, \ldots, x_{n}) \in I(X)\).
                Hence \(F^{n}(P) = 0, \forall P \in X\).
                But note that \(F^{n}(P) =(F(P))^{n} \implies (F(P))^{n} = 0 \implies F \in I(X)\).

            \item We can simply apply the criterion for a function to be one-to-one.
                In fact, if \(X_{1}, X_{2}\) are closed affine sets such that \(I(X_{1}) = I(X_{2})\),
                then trivially \(V(I(X_{1})) = V(I(X_{2})) \implies X_{1} = X_{2}\).\qed
        \end{enumerate}
    \end{proof*}

    \subsection{Noetherian topological spaces}
    \subfile{../sub/2.1 - Noetherian topological spaces}


\end{document}
